\section{Encoder, decoder y las tres familias de modelos}

El Transformer original se diseñó para la traducción automática y consta de un encoder y un decoder. Ambos apilan bloques de attention y FFN, pero difieren en el acceso a la información. El \term{encoder} (codificador) suele permitir que cada token de origen lea la secuencia de origen en ambas direcciones; el \term{decoder} (decodificador), al generar el $t$-ésimo token de destino, solo puede leer el prefijo de destino y, mediante cross-attention, lee las encoder outputs.

\subsection{Flujo de datos encoder-decoder}

Durante el entrenamiento, los source tokens pasan por el encoder y producen la memory; los target tokens, desplazados a la derecha, entran en el masked decoder. La self-attention del decoder procesa el prefijo ya generado; la \term{cross-attention} usa el decoder state como query y la encoder memory como key/value, y al final una softmax predice el siguiente token. \term{autoregressive} (autorregresivo) significa que la probabilidad conjunta se descompone como
\[
p(y\mid x)=\prod_{t=1}^{m}p(y_t\mid y_{<t},x).
\]

\lecturefigure{slide-14-encoder-decoder.jpg}{El Transformer original realiza seq2seq con la encoder memory, el masked decoder y la cross-attention.}{Video oficial 13:40; Vaswani et al. 2017.}

\begin{knowledgebox}{Lectura de la figura: no hay que mezclar los tres tipos de attention}
La self-attention del encoder lee todo el source; la masked self-attention del decoder solo lee el target prefix; la cross-attention permite que la target query lea la encoder memory. En la figura, el target desplazado a la derecha garantiza que el modelo no pueda ver directamente la respuesta actual. El entrenamiento puede calcular en paralelo todas las target positions, pero la máscara de cada position conserva la semántica autorregresiva; la inferencia sigue requiriendo avanzar token por token o recurrir a métodos especializados de decodificación en paralelo.
\end{knowledgebox}

\begin{importantbox}{Las tres familias principales}
\begin{enumerate}
  \item \textbf{Encoder-only}: representación bidireccional, adecuada para clasificación, recuperación, etiquetado y representation learning; ejemplo representativo: BERT.
  \item \textbf{Decoder-only}: causal next-token modeling, adecuada para la generación abierta y para una prompt interface unificada; ejemplo representativo: la serie GPT.
  \item \textbf{Encoder-decoder}: source y target separados, adecuada para traducción, resumen y generación condicionada; ejemplos representativos: el Transformer original y T5.
\end{enumerate}
Elegir una familia es, ante todo, elegir los límites de la información y la interfaz de la tarea, no solo la cantidad de parámetros.
\end{importantbox}

\subsection{Resumen de la sección}

El encoder, el decoder y la cross-attention definen de dónde proviene la información y cuándo es visible. La diferencia entre GPT y BERT puede deducirse directamente de este diagrama de la arquitectura original, en lugar de memorizarse como dos nombres de modelos aislados.
